%==============================================================================
% main.tex — template for mathematical work
% Styles inherited from last_paper.tex
% Build:  latexmk -pdf -outdir=build main.tex
%==============================================================================
\documentclass{article}

% Language setting
% Replace `english' with e.g. `spanish' to change the document language
\usepackage[english]{babel}
\usepackage[colorinlistoftodos]{todonotes}

% Set page size and margins
\usepackage[a4paper,top=2cm,bottom=2cm,left=3cm,right=3cm,marginparwidth=1.75cm]{geometry} % old
% \usepackage[a4paper, top=1.5cm, bottom=1.5cm, left=1.5cm, right=1.5cm, marginparwidth=1.2cm]{geometry} % very small paddings

% Useful packages
\usepackage{amsmath, amsthm}
\usepackage{amssymb}
\usepackage{graphicx}
\usepackage[hidelinks]{hyperref}
\usepackage{xcolor}
\usepackage{thmtools}
\usepackage{float}
\usepackage[framemethod=TikZ]{mdframed}

\graphicspath{{figures/}}

%------------------------------------------------------------------------------
% Extras beyond the inherited preamble — delete what you do not need
%------------------------------------------------------------------------------
\usepackage{mathtools}   % \DeclarePairedDelimiter, \coloneqq, amsmath fixes
\usepackage{booktabs}    % \toprule / \midrule / \bottomrule
\usepackage{enumitem}    % list spacing control
\usepackage{listings}
\lstset{
    basicstyle=\ttfamily\small,
    breaklines=true,
    frame=single,
    numbers=left,
    numberstyle=\tiny,
    showstringspaces=false,
}

%------------------------------------------------------------------------------
% Macros
%------------------------------------------------------------------------------
% Empirical estimators
\newcommand{\Covhat}{\widehat{\mathrm{Cov}}}
\newcommand{\Varhat}{\widehat{\mathrm{Var}}}

% Number sets (optional shorthands — the paper style writes \mathbb{E} inline)
\newcommand{\R}{\mathbb{R}}
\newcommand{\N}{\mathbb{N}}
\newcommand{\Z}{\mathbb{Z}}
\newcommand{\Q}{\mathbb{Q}}
\newcommand{\C}{\mathbb{C}}
\newcommand{\E}{\mathbb{E}}
\newcommand{\Prob}{\mathbb{P}}
\newcommand{\eps}{\varepsilon}

% Auto-sizing delimiters: \abs{x} fixed size, \abs*{x} scales to content
\DeclarePairedDelimiter{\abs}{\lvert}{\rvert}
\DeclarePairedDelimiter{\norm}{\lVert}{\rVert}
\DeclarePairedDelimiter{\inner}{\langle}{\rangle}
\DeclarePairedDelimiter{\set}{\{}{\}}
\DeclarePairedDelimiter{\ceil}{\lceil}{\rceil}
\DeclarePairedDelimiter{\floor}{\lfloor}{\rfloor}

% Operators
\DeclareMathOperator*{\argmin}{arg\,min}
\DeclareMathOperator*{\argmax}{arg\,max}
\DeclareMathOperator{\tr}{tr}
\DeclareMathOperator{\rank}{rank}
\DeclareMathOperator{\diag}{diag}
\DeclareMathOperator{\spann}{span}

%------------------------------------------------------------------------------
% Theorem environments
%------------------------------------------------------------------------------
% new style for definition
\declaretheoremstyle[
    mdframed={
        backgroundcolor=gray!10, 
        linecolor=black, 
        linewidth=0.5pt,
        innertopmargin=10pt,
        innerbottommargin=10pt,
        nobreak=true
    },
    headfont=\bfseries,
    notefont=\bfseries, 
    notebraces={ (}{)},
]{shadedstyle}

% plain (unshaded) companion style for the supporting statements
\declaretheoremstyle[
    headfont=\bfseries,
    notefont=\bfseries,
    notebraces={ (}{)},
    bodyfont=\itshape,
]{plainstyle}

\declaretheorem[style=shadedstyle, name=Definition]{definition}
\declaretheorem[style=shadedstyle, name=Theorem]{theorem}
\declaretheorem[style=plainstyle, name=Lemma, sibling=theorem]{lemma}
\declaretheorem[style=plainstyle, name=Proposition, sibling=theorem]{proposition}
\declaretheorem[style=plainstyle, name=Corollary, sibling=theorem]{corollary}
\declaretheorem[style=remark, name=Remark, sibling=theorem]{remark}
\declaretheorem[style=definition, name=Example, sibling=theorem]{example}

%------------------------------------------------------------------------------
\title{Title of the Work}
\author{Aleksandr Mazovetskii}

\begin{document}

\begin{titlepage}
    \vspace*{\fill} 
    \begin{center}
        {\Huge \bfseries Title of the Work\par}
        \vspace{1.5cm}
        {\Large Aleksandr Mazovetskii\par}
        \vspace{1cm}
        {\large \today\par}
    \end{center}
    \vspace*{\fill} 
    
\end{titlepage}

\tableofcontents
\newpage

\section{Introduction}
One paragraph on the problem, why the existing approach is insufficient, and what
this work establishes. Close with a sentence naming the main application.

\section{Preliminaries}
\subsection{Setting}
Fix the objects and notation used throughout. Reference statements the way the
paper does, with Definition~\ref{def:example} and Theorem~\ref{thm:main}.

\begin{definition}[Name of the notion] \label{def:example}
Let $X$ be a set. A function $f: X \to \mathbb{R}$ is called \textbf{bounded} if
there exists $M \in \mathbb{R}_{\ge 0}$ such that
$$
\forall x \in X: \quad |f(x)| \le M
$$
\textbf{Note:} remarks that belong inside the shaded box go here.
\end{definition}

Display equations carry a label only when they are referenced later:
\begin{equation} \label{eq:example}
\int_{a}^{b} f'(x)\,dx = f(b) - f(a).
\end{equation}

Multi-line derivations align on the relation symbol, with the justification of
each step set in text:
\begin{align*}
\|x + y\|^2 &= \langle x + y,\, x + y \rangle \\
&= \|x\|^2 + 2\langle x, y \rangle + \|y\|^2
\hspace{0.5cm} \text{(by bilinearity)}
\end{align*}

\section{Main results}

\begin{theorem}[Main theorem] \label{thm:main} \cite{example_reference}
The statement to be proved, under the hypotheses of Definition~\ref{def:example}.
\end{theorem}

\begin{proof}
The argument. Equation~\eqref{eq:example} is used here.
\end{proof}

\begin{lemma} \label{lem:auxiliary}
An auxiliary statement used in the proof above.
\end{lemma}

\begin{proof}
Argument.
\end{proof}

\begin{remark}
A comment that sits outside the formal development.
\end{remark}

\section{Experiments}

\begin{figure}[H]
    \centering
    % \includegraphics[width=0.75\linewidth]{example_plot.png}
    \caption{What the reader should take from the plot, followed by the
    parameters used: $N=10\,000, M=100$}
    \label{fig:example}
\end{figure}

\begin{table}[H]
    \centering
    \begin{tabular}{lrr}
        \toprule
        Method & Error & Time (s) \\
        \midrule
        Baseline & 0.42 & 1.3 \\
        Proposed & 0.17 & 2.1 \\
        \bottomrule
    \end{tabular}
    \caption{Caption.}
    \label{tab:example}
\end{table}

\section{Conclusion}
What was shown, what the experiments confirm, and the natural next step.

\appendix
\section{Deferred proofs}
Longer computations that would interrupt the main text.

\newpage
\bibliographystyle{plain}
\bibliography{literature}

\end{document}
